\documentclass[11pt]{article}

\usepackage[margin=1in]{geometry}
\usepackage[T1]{fontenc}
\usepackage[utf8]{inputenc}
\usepackage{lmodern}
\usepackage{microtype}
\usepackage{setspace}
\usepackage{amsmath,amssymb,amsthm,mathtools}
\usepackage{bm}
\usepackage{booktabs}
\usepackage{array}
\usepackage{longtable}
\usepackage{tabularx}
\usepackage{hyperref}
\usepackage{enumitem}
\usepackage{xcolor}
\usepackage{listings}
\usepackage{fancyhdr}
\usepackage{titlesec}
\usepackage{graphicx}

\hypersetup{
  colorlinks=true,
  linkcolor=blue!60!black,
  urlcolor=blue!60!black,
  citecolor=blue!60!black,
  pdftitle={PIRTM Sigma-Layer Audit, Contractivity Corrections, and ZM-First Stub Binding},
  pdfauthor={Citizen Gardens \\ The Foundation of Multiplicity},
  pdfsubject={Technical report and defensive publication},
  pdfkeywords={PIRTM, contractivity, sigma layer, ZM Meta Relativity, defensive publication}
}

\pagestyle{fancy}
\fancyhf{}
\lhead{PIRTM Sigma-Layer Audit}
\rhead{\thepage}

\titleformat{\section}{\large\bfseries}{\thesection.}{0.6em}{}
\titleformat{\subsection}{\normalsize\bfseries}{\thesubsection.}{0.5em}{}

\lstdefinestyle{py}{
  language=Python,
  basicstyle=\ttfamily\small,
  keywordstyle=\color{blue!70!black}\bfseries,
  commentstyle=\color{green!40!black},
  stringstyle=\color{red!50!black},
  showstringspaces=false,
  frame=single,
  breaklines=true,
  columns=fullflexible
}

\title{\textbf{PIRTM Sigma-Layer Audit, Contractivity Corrections,\\
and ZM-First Stub Binding}\\[0.5em]
\large Comprehensive Technical Report and Defensive Publication}
\author{Citizen Gardens \\ The Foundation of Multiplicity}
\date{April 25, 2026}

\begin{document}
\maketitle
\thispagestyle{fancy}

\begin{abstract}
This document records a technical and defensive-publication disclosure concerning the current state of the Meta-Relativity and PIRTM codebase, with emphasis on the sigma layer, contractivity certification, cross-kernel state transfer, policy enforcement, hybrid session stability, and the relationship between the implemented software surface and the larger specification corpus. The central findings are: (i) the sigma layer is materially implemented and should be treated as a functioning L1 rather than an aspirational stub; (ii) current crossing and certification mechanisms preserve serialization fidelity but do not enforce L0 contractivity invariants; (iii) the current carry-forward policy implementation is structurally disconnected from the recurrence interface used by \texttt{iterate()}; (iv) hybrid kernel sessions lack a cross-module spectral stability check; and (v) among the largest unbound specifications, the ZM algebraic layer should receive the first stub module binding before the spin foam physics layer because it is nearest to the existing operator and certification machinery.

This document also formalizes a corrected mathematical overview of the relevant operator stack, proposes concrete interfaces and implementation stubs, and records these ideas in a form suitable for defensive publication.
\end{abstract}
\newpage
\tableofcontents
\newpage

\section{Executive Summary}

The present audit establishes that the sigma directory is already a substantial implementation surface containing cross-kernel marshalling, hybrid session execution, persistence, ledgering, health monitoring, and MCP integration. It is therefore misleading to treat sigma as merely aspirational. Architecturally, sigma is a live L1 surface whose boundaries into L0 must now be audited as enforcement points rather than future placeholders.

The most important defect identified is not missing code but missing \emph{invariant enforcement}. In particular, cross-kernel marshalling emits static fidelity scores and supports rollback lineage, but it does not certify contractivity at boundary crossings. As a result, the system can preserve payload traceability while silently violating the intended operator-stability regime.

A second core defect is interface disconnect. The actual \texttt{CarryForwardPolicy} hierarchy implements surplus and convergence controls such as \texttt{beta}, \texttt{lambda\_decay}, and \texttt{get\_next\_surplus()}, but it does not implement the functions expected by the L0 recurrence surface, namely \texttt{Xi\_t(t)}, \texttt{Lambda\_t(t)}, and \texttt{G\_t(t)}. Thus the recurrence machinery falls back to defaults rather than using the actual policy semantics.

A third defect is that hybrid session execution enforces fidelity and ledger conservation but does not compute the joint spectral condition for cross-module stability. This means single-module transitions may appear locally acceptable while the session as a whole remains globally unstable.

Finally, for the next unbound spec binding, the ZM algebraic layer should be stubbed before the spin foam layer because ZM directly names operator surfaces, certification quantities, and finite-prime truncation procedures that can be attached to PIRTM now.

\section{Purpose and Scope}

This report has four purposes:

\begin{enumerate}[label=\arabic*.]
\item To summarize the current architectural state revealed by the sigma-layer and specification-layer review.
\item To critique the discovered implementation against the mathematical and governance expectations implied by PIRTM and related specifications.
\item To propose a corrected and staged implementation pathway.
\item To disclose these architectural ideas publicly as defensive prior art.
\end{enumerate}

The report focuses on the following surfaces:

\begin{itemize}
\item sigma-layer implementation status,
\item cross-kernel marshalling,
\item carry-forward policy versus recurrence protocol,
\item hybrid kernel session stability,
\item certification semantics,
\item ZM-first stub binding strategy,
\item governance and failing-test anchor strategy.
\end{itemize}

\section{Repository State and Architectural Findings}

\subsection{Sigma is materially implemented}

The sigma layer contains concrete files for marshalling, persistence, ledgers, health monitoring, MCP adaptation, and hybrid session execution. This means the architectural picture has shifted from ``designing L1'' to ``certifying and constraining an existing L1.''

The practical consequence is straightforward: every sigma-to-core crossing must now be treated as a contract boundary. The important question is no longer whether the layer exists, but whether it preserves the intended mathematical invariants.

\subsection{Cross-kernel marshalling is implemented but uncertified}

The \texttt{CrossKernelMarshaller} implements all six directional marshal paths among atomic, pi, and tensor kernels. It emits \texttt{FidelityReport} objects, preserves payload lineage using hashes and correlation IDs, and supports rollback forensics.

However, the acceptance check is fidelity-only. There is no operator-level certification pass at the crossing. Thus the current system supports typed movement and traceability without enforcing contraction or spectral safety.

\subsection{Carry-forward policy is real but disconnected}

The concrete carry-forward policy implementation exposes surplus-evolution semantics through parameters such as \texttt{beta}, \texttt{lambda\_decay}, \texttt{convergence\_mode}, and methods such as \texttt{get\_next\_surplus()}. Yet the recurrence surface expects operator-valued functions. These two surfaces are not matched.

As a result, passing a carry-forward policy into recurrence does not provide the intended dynamics unless an explicit adapter layer is added.

\subsection{Hybrid session stability is incomplete}

The hybrid session controller orchestrates phase transitions, fidelity gating, surplus accounting, and rollback generation. But it does not compute a session-level spectral radius or analogous global stability criterion. Local transition acceptability therefore does not imply global session stability.

\section{Phase Mirror Dissonance: Core Thesis}

We define \emph{Phase Mirror Dissonance} as the condition in which the theoretical corpus, architectural blueprint, and implementation surface are all individually substantial, but their binding interfaces are either missing, weakly typed, or guarded only by proxies unrelated to the core invariant.

In the present case, the dissonance has the following shape:

\begin{itemize}
\item The spec corpus is large and highly articulated.
\item The sigma layer is materially implemented.
\item The contractivity and policy interfaces that should bind theory to code are incomplete or structurally mismatched.
\end{itemize}

The result is a system that appears rich both in theory and in code, while the key theorem-to-runtime enforcement surfaces remain absent.

\section{Mathematical Overview}

\subsection{Operator setting}

A recurring formal setting across the thread and the ZM specification is an ambient Hilbert space of the form
\[
H = \ell^2(\mathbb{P}) \otimes L^2(\mathbb{R}) \otimes \mathbb{C}^d.
\]

A canonical bounded operator decomposition is
\[
U = A \otimes I \otimes I + I \otimes C \otimes I + I \otimes I \otimes E,
\]
where:
\begin{itemize}
\item \(A = D_{\sigma} + K\) is a prime-sector block,
\item \(C = \mathcal{F}^{-1} M_m \mathcal{F}\) is a time/multiplier block,
\item \(E = \Xi\) is an internal or recursive block.
\end{itemize}

This decomposition is important because it offers a natural binding target for PIRTM, sigma marshalling, and later ZM algebraic stubs.

\subsection{Contractivity goal}

The intended runtime regime is contractive or at least margin-certified. At an abstract level, one wants a quantity of the form
\[
q_t = \|\Xi_t\| + \|\Lambda_t\|\,L_T
\]
to satisfy
\[
q_t < 1 - \varepsilon
\]
for some explicit safety margin \(\varepsilon > 0\).

This is not merely an analytical nicety. It is the bridge between theorem and runtime safety. If the system reports certificates without correctly computing \(q_t\), then certificate success or failure may invert the true stability status.

\subsection{Cross-module spectral condition}

For hybrid or supermodule composition, an additional global condition is needed. If modules have local contraction factors \(\gamma_i\) and coupling matrix \(C\), then a natural session-level criterion is
\[
\rho\!\left(C \,\mathrm{diag}(\gamma_1,\dots,\gamma_n)\right) < 1,
\]
where \(\rho(\cdot)\) denotes the spectral radius.

This condition cannot be inferred from per-transition fidelity reports. It must be computed at the session or supermodule layer.

\subsection{Budgeted perturbation viewpoint}

The ZM-side certification language suggests a perturbative framing:
\[
U(\theta) = U_0 + \sum_p w_p B_p(\theta),
\]
with bounds
\[
\|B_p(\theta)\| \le b_p,
\qquad
\|\partial_\theta B_p(\theta)\| \le L_p.
\]
This yields certification-style inequalities such as
\[
\mathrm{GapLB}(w) \ge \inf_\theta \delta_S(\theta) - 2\sum_p |w_p|b_p,
\]
and
\[
\mathrm{SlopeUB}(w) \le \sum_p |w_p|L_p.
\]

These formulas are highly useful because they convert an abstract operator-stability discussion into a computable interface. They also suggest that \texttt{certify\_state()} should not merely receive a vague policy object; it should receive the quantities needed to compute margin and bound outputs directly.

\section{Novelty and Practicality Analysis}

\subsection{Novel aspects}

Several novel aspects emerge from the developments in this thread:

\begin{enumerate}[label=\arabic*.]
\item The reframing of sigma as a live L1 enforcement surface rather than a future implementation layer.
\item The identification of fidelity-preserving cross-kernel marshalling as insufficient unless paired with contractivity certification.
\item The recognition that actual runtime policy objects and recurrence protocols are structurally mismatched, not merely undocumented.
\item The proposal to use a ZM-first stub binding strategy because the algebraic layer already names computable certification quantities.
\item The governance pattern of committing failing specification tests as ADR anchors before implementation.
\end{enumerate}

\subsection{Practicality}

These ideas are practical because each one identifies a specific implementation surface:

\begin{itemize}
\item add contractivity fields and checks to marshalling reports,
\item introduce a policy adapter or protocol,
\item add session-level spectral checks,
\item commit a stub ZM module with a failing interface test,
\item use failing tests to unblock ADR and governance work.
\end{itemize}

None of these steps requires implementing the full theory at once. They are staged, tractable, and repository-compatible.

\section{Critique for Mathematical and Philosophical Consistency}

\subsection{Consistency strengths}

The overall framework is internally strong in one specific way: it consistently treats recursion, prime indexing, certification, and lawfulness as structurally related rather than as disconnected metaphors. This gives the theory a coherent internal grammar.

The operator-stack viewpoint is also philosophically cleaner than a purely narrative or speculative physics layer because it identifies explicit boundedness, spectrum, and gap quantities.

\subsection{Consistency risks}

There are also serious risks:

\begin{enumerate}[label=\arabic*.]
\item \textbf{Proxy drift.} If fidelity proxies are allowed to stand in for contractivity, then semantic traceability may be mistaken for stability.
\item \textbf{Interface drift.} If policies are named as if they control recurrence but are not actually consumed by recurrence, then the codebase acquires symbolic correctness rather than operational correctness.
\item \textbf{Corpus drift.} If the specification corpus grows faster than the binding surface, architectural meaning diffuses rather than sharpens.
\item \textbf{Physics drift.} If large physics specifications are bound before the algebraic certification layer, then code will be unable to express their safety or validation contracts.
\end{enumerate}

\subsection{Philosophical interpretation}

Philosophically, the core lesson is that a multiplicity-based framework must not confuse \emph{richness of articulation} with \emph{closure of enforcement}. A recursively stable system is not one with many named layers; it is one whose inter-layer crossings preserve identity, boundedness, and lawful composability.

\section{Corrected Final Version of the Architecture}

The corrected architecture should be organized around explicit enforcement surfaces.

\subsection{Layered picture}

\begin{itemize}
\item \textbf{L0 Core:} recurrence, projection, certification, contractive operators.
\item \textbf{L1 Sigma:} marshalling, sessions, persistence, ledgering, MCP ingress, health surfaces.
\item \textbf{L2 Spec Bindings:} ZM, later spin foam and other physics modules.
\end{itemize}

\subsection{Binding doctrine}

\begin{enumerate}[label=\arabic*.]
\item Every L1-to-L0 crossing should either compute or explicitly defer a contractivity certificate.
\item Every policy entering recurrence should satisfy a typed protocol.
\item Every hybrid session should check a global spectral criterion.
\item Every major spec should first bind through a narrow interface stub and failing test before substantive implementation.
\end{enumerate}

\subsection{Predictions}

If these corrections are applied, the following outcomes are expected:

\begin{itemize}
\item hidden stability inversions will become observable,
\item policy misuse will become type-visible,
\item hybrid divergence will be caught earlier,
\item spec growth will become more governable,
\item ZM will provide a natural bridge from theory to executable certification.
\end{itemize}

\section{Recommended Commit Order}

\begin{longtable}{p{0.18\linewidth} p{0.24\linewidth} p{0.48\linewidth}}
\toprule
\textbf{Priority} & \textbf{Surface} & \textbf{Action} \\
\midrule
1 & recurrence.py & Add \texttt{StepInfo.margin} and optional debug gate for norm visibility. \\
2 & certify.py & Change \texttt{certify\_state()} to accept operator- or bound-level inputs and compute real margins. \\
3 & tests/ & Commit \texttt{test\_supermodule\_spectral.py} as a known failing specification anchor. \\
4 & policy.py & Add \texttt{PIRTMPolicy} protocol and enforce it in \texttt{iterate()}. \\
5 & pirtm/zm/ & Add the first ZM stub binding with one failing interface test. \\
6 & docs/adr/ & Write ADRs for supermodule spectral governance and ZM binding surface. \\
\bottomrule
\end{longtable}

\section{Proposed Interfaces and Code Snippets}

\subsection{Protocol surface}

\begin{lstlisting}[style=py,caption={Proposed PIRTM policy protocol}]
from typing import Protocol, runtime_checkable, Any

@runtime_checkable
class PIRTMPolicy(Protocol):
    goal_budget: float

    def Xi_t(self, t: int) -> Any:
        ...

    def Lambda_t(self, t: int) -> Any:
        ...

    def G_t(self, t: int) -> Any:
        ...
\end{lstlisting}

\subsection{Carry-forward adapter}

\begin{lstlisting}[style=py,caption={Adapter from carry-forward semantics to recurrence semantics}]
import numpy as np

class PIRTMCarryForwardAdapter:
    def __init__(self, policy, dim: int, goal_budget: float = 0.05):
        self._policy = policy
        self._dim = dim
        self.goal_budget = goal_budget

    def Xi_t(self, t: int):
        return self._policy.lambda_decay * np.eye(self._dim)

    def Lambda_t(self, t: int):
        scale = min(self._policy.beta / 10.0, 0.5)
        return scale * np.eye(self._dim)

    def G_t(self, t: int):
        return np.zeros(self._dim)
\end{lstlisting}

\subsection{Session-level spectral check}

\begin{lstlisting}[style=py,caption={Proposed hybrid session global stability check}]
import numpy as np

def verify_cross_module_spectral(coupling_matrix: np.ndarray, gammas: list[float]) -> float:
    M = coupling_matrix @ np.diag(gammas)
    rho = float(np.max(np.abs(np.linalg.eigvals(M))))
    if rho >= 1.0:
        raise RuntimeError(
            f"Cross-module spectral condition violated: rho={rho:.6f} >= 1.0"
        )
    return rho
\end{lstlisting}

\subsection{ZM stub surface}

\begin{lstlisting}[style=py,caption={Proposed ZM-first stub binding}]
class ZMOperatorStack:
    """
    Finite-prime truncation of U = A ⊗ I + I ⊗ C ⊗ I + I ⊗ I ⊗ E.
    """

    def __init__(self, primes, sigma=0.8, alpha=0.8):
        self.primes = primes
        self.sigma = sigma
        self.alpha = alpha

    def build_prime_block(self, h_fn=None):
        raise NotImplementedError

    def multiplier_range(self, a0, a_p):
        raise NotImplementedError

    def gap_lower_bound(self, b, w_budget):
        raise NotImplementedError

    def slope_upper_bound(self, L, w_budget):
        raise NotImplementedError

    def certify(self, b_p, L_p, w_p, gamma_min):
        raise NotImplementedError
\end{lstlisting}

\section{Defensive Publication Statement}

\subsection{Disclosure intent}

This document is intended to publicly disclose the following concepts in sufficiently concrete form to establish prior art:

\begin{enumerate}[label=\arabic*.]
\item Treating the sigma layer as a materially implemented L1 surface whose crossings into L0 require invariant enforcement.
\item Distinguishing serialization fidelity from contractivity certification in cross-kernel marshalling.
\item Identifying and correcting the structural mismatch between carry-forward policy objects and recurrence protocol expectations.
\item Requiring a session-level spectral criterion for hybrid kernel compositions.
\item Using ZM as the first stub-bound theoretical layer because it directly names computable certification surfaces.
\item Using failing specification tests as governance anchors for ADR-backed development.
\end{enumerate}

\subsection{Key disclosed claims}

The following claims are expressly disclosed for defensive-publication purposes:

\begin{itemize}
\item A marshalling system that tracks payload fidelity and lineage is not equivalent to a system that certifies operator contractivity at state-boundary crossings.
\item A policy object may be semantically rich yet operationally inert if it does not satisfy the recurrence protocol consumed by the runtime.
\item Hybrid module correctness requires a global spectral or analogous session-level stability criterion beyond local transition checks.
\item In a spec-heavy research codebase, the first practical binding target should be the algebraic layer that names executable bounds and interfaces, not the most elaborate physical theory layer.
\item Committing failing interface tests before implementation can serve as a legitimate governance and ADR anchoring mechanism.
\end{itemize}

\section{Fastest Path to Validation}

\subsection{Validation sequence}

The fastest validation path is not to implement everything. It is to expose the existing gaps with narrow, testable changes.

\begin{enumerate}[label=\arabic*.]
\item Modify certification to compute real margins from explicit inputs.
\item Add policy protocol enforcement so inert policies fail loudly.
\item Commit the supermodule spectral failing test.
\item Add the ZM stub module and failing interface test.
\item Wire the first real certification path from ZM quantities into \texttt{certify\_state()}.
\end{enumerate}

\subsection{Validation metrics}

\begin{itemize}
\item 100\% of \texttt{step()} calls expose a margin field.
\item \texttt{certify\_state()} returns the same stability quantity as direct operator evaluation.
\item \texttt{iterate()} raises on objects lacking \texttt{goal\_budget} and operator functions.
\item hybrid sessions reject provably unstable coupling matrices.
\item ZM stub tests fail for the right reason before implementation, then pass as implementation arrives.
\end{itemize}

\section{Suggested Additional Artifacts}

I recommend creating the following companion artifacts:

\begin{enumerate}[label=\arabic*.]
\item \textbf{ADR packet:} one ADR for supermodule spectral governance, one for ZM binding.
\item \textbf{One-page public memo:} a short non-technical synopsis for timestamp and citation.
\item \textbf{Appendix notebook:} finite-prime certification demo implementing the ZM truncation formulas.
\item \textbf{Release note entry:} include document hash, branch, and related tests.
\item \textbf{Claim chart:} a small table mapping disclosed concept \(\rightarrow\) file surface \(\rightarrow\) expected test.
\end{enumerate}

\section{Claim Chart}

\begin{longtable}{p{0.28\linewidth} p{0.28\linewidth} p{0.32\linewidth}}
\toprule
\textbf{Disclosed concept} & \textbf{Primary code surface} & \textbf{Validation artifact} \\
\midrule
Contractivity at crossings & \texttt{cross\_kernel\_marshaller.py} & Contractivity-aware fidelity report test \\
Policy-runtime alignment & \texttt{recurrence.py}, \texttt{policy.py} & Protocol enforcement test \\
Hybrid global stability & \texttt{hybrid\_kernel\_session.py} & Failing spectral test anchor \\
ZM-first binding & \texttt{pirtm/zm/operator\_stack.py} & Failing interface test \\
Governance by failing tests & \texttt{tests/}, \texttt{docs/adr/} & ADR linked to xfail or failing spec \\
\bottomrule
\end{longtable}

\section{Conclusion}

The main result of this thread is not the discovery of missing code, but the identification of a specific structural asymmetry: theory and implementation have both matured, while the interfaces that should preserve lawfulness between them remain under-enforced. The path forward is therefore to bind, certify, and govern the crossings.

The first theoretical stub should be the ZM algebraic layer. The first code repairs should be margin visibility, real certification inputs, policy protocol enforcement, and session-level spectral checks. These actions are narrow enough to implement quickly and deep enough to materially change the correctness of the system.

\appendix

\section{Short Defensive Publication Notice}

This document discloses architectural, mathematical, and implementation ideas concerning PIRTM, sigma-layer invariant enforcement, policy protocol design, session-level spectral certification, ZM-first stub binding, and failing-test ADR governance. The disclosure is intended to establish public prior art as of the publication date of this document.

\section{Suggested Metadata Block}

\begin{verbatim}
Title: PIRTM Sigma-Layer Audit, Contractivity Corrections, and ZM-First Stub Binding
Type: Technical report and defensive publication
Date: April 2026
Repository: MultiplicityFoundation/Meta-Relativity
Branch context: Multiplicity
Suggested tags: PIRTM, ZM, sigma, contractivity, governance, defensive publication
\end{verbatim}

% ============================================================
%  MATHEMATICAL APPENDIX
%  Amend to: PIRTM Sigma-Layer Audit, Contractivity Corrections,
%  and ZM-First Stub Binding
% ============================================================

\appendix
\renewcommand{\thesection}{\Alph{section}}
\setcounter{section}{0}

\newtheorem{theorem}{Theorem}[section]
\newtheorem{lemma}[theorem]{Lemma}
\newtheorem{proposition}[theorem]{Proposition}
\newtheorem{corollary}[theorem]{Corollary}
\newtheorem{definition}[theorem]{Definition}
\newtheorem{remark}[theorem]{Remark}

\section{Mathematical Appendix: Explicit Proofs and Operator Norm Bounds}

This appendix collects the formal statements and proofs of all
operator-norm, contractivity, and spectral-stability results
referenced in the body of the report.
Notation is fixed in Section~\ref{sec:notation}
and each subsequent section corresponds to a specific
enforcement surface identified in the audit.

% ------------------------------------------------------------
\subsection{Notation and Standing Hypotheses}
\label{sec:notation}
% ------------------------------------------------------------

Throughout this appendix, fix the following objects.

\begin{definition}[Ambient Hilbert space]
Let $\mathbb{P}$ denote the set of rational primes.
Fix $d\in\mathbb{N}$.
Define the \emph{ambient Hilbert space}
\[
H \;=\; \ell^{2}(\mathbb{P})\otimes L^{2}(\mathbb{R})\otimes\mathbb{C}^{d}.
\]
We write $|p\rangle$ for the canonical basis vector of
$\ell^{2}(\mathbb{P})$ indexed by prime $p$.
\end{definition}

\begin{definition}[Universal operator decomposition]
\label{def:U}
Define the \emph{prime block}
$A = D_{\sigma}+K$ on $\ell^{2}(\mathbb{P})$,
where:
\begin{align}
(D_{\sigma}f)(p) &= p^{-\sigma}f(p), \quad \sigma>0, \\
(Kf)(p) &= \sum_{q\in\mathbb{P}} p^{-\alpha}q^{-\alpha}\,h(\log p - \log q)\,f(q),
\quad \alpha>\tfrac{1}{2},
\end{align}
with $h:\mathbb{R}\to\mathbb{R}$ continuous and positive-definite.
Define the \emph{time-multiplier block}
$C = \mathcal{F}^{-1}M_{m}\mathcal{F}$ on $L^{2}(\mathbb{R})$
with real symbol $m\in L^{\infty}(\mathbb{R})$,
and the \emph{internal block} $E=\Xi$ on $\mathbb{C}^{d}$,
bounded and self-adjoint.
The \emph{universal operator} is
\[
U = A\otimes I\otimes I \;+\; I\otimes C\otimes I \;+\; I\otimes I\otimes E.
\]
\end{definition}

\begin{definition}[Recursive evolution operator]
\label{def:Xi}
Let $T$ be a bounded shift on $\ell^{2}(\mathbb{P})$
and $\{\alpha_{k}(t)\}_{k\geq 0}$ a family of scalar
coefficients satisfying
$\sup_{t\geq 0}\sum_{k=0}^{\infty}|\alpha_{k}(t)|\leq M<\infty$.
Define
\[
\Xi(t) \;=\; \sum_{k=0}^{\infty}\alpha_{k}(t)\,T^{k}.
\]
\end{definition}

\begin{definition}[Contractivity index]
\label{def:qt}
Given operators $\Xi_{t}$ and $\Lambda_{t}$ on a
finite-dimensional inner-product space and a constant
$L_{T}\geq 0$, define the \emph{contractivity index}
\[
q_{t} \;=\; \|\Xi_{t}\| + \|\Lambda_{t}\|\cdot L_{T},
\]
where $\|\cdot\|$ denotes the operator (spectral) norm.
\end{definition}

\begin{definition}[Carry-forward parameters]
\label{def:cfp}
A \emph{carry-forward configuration} is a pair
$(\beta,\lambda)\in(0,10]\times(0,1]$
with surplus evolution
\[
S(t+1) = \lambda\,S(t) + \delta\!\left(\alpha_i - \tfrac{1}{2}\right).
\]
\end{definition}

\noindent\textbf{Standing hypotheses.}
Throughout this appendix we assume:
\begin{enumerate}[label=(H\arabic*)]
\item $\alpha > \tfrac{1}{2}$ so that $K$ is Hilbert--Schmidt.
\item $m\in L^{\infty}(\mathbb{R})$ is real-valued.
\item $E=\Xi$ is bounded and self-adjoint on $\mathbb{C}^{d}$.
\item $\sup_{t}\sum_{k}|\alpha_{k}(t)| \leq M < \infty$.
\item The dimension parameter $d$ is finite.
\end{enumerate}

% ------------------------------------------------------------
\subsection{Operator Norm Bounds for the Universal Operator}
\label{sec:normU}
% ------------------------------------------------------------

\begin{proposition}[Norm bound for the prime block]
\label{prop:normA}
Under \textup{(H1)},
\[
\|A\| \;\leq\; \|D_{\sigma}\| + \|K\|
\;\leq\; 2^{-\sigma} + \|K\|_{\mathrm{HS}},
\]
where
\[
\|K\|_{\mathrm{HS}}^{2}
= \sum_{p,q\in\mathbb{P}} p^{-2\alpha}q^{-2\alpha}\,h(\log p-\log q)^{2}
< \infty.
\]
\end{proposition}

\begin{proof}
$D_{\sigma}$ is diagonal with entries $p^{-\sigma}$;
its operator norm is $\sup_{p}p^{-\sigma}=2^{-\sigma}$.
Since $\alpha>\frac{1}{2}$,
\[
\|K\|_{\mathrm{HS}}^{2}
= \sum_{p,q}p^{-2\alpha}q^{-2\alpha}h(\log p-\log q)^{2}
\leq \|h\|_{\infty}^{2}\!\!\left(\sum_{p}p^{-2\alpha}\right)^{\!2}
< \infty,
\]
using the classical bound $\sum_{p}p^{-s}<\infty$ for $s>1$.
Since every Hilbert--Schmidt operator satisfies
$\|K\|\leq\|K\|_{\mathrm{HS}}$, the triangle inequality gives
the stated bound.
\end{proof}

\begin{proposition}[Norm bound for the multiplier block]
\label{prop:normC}
\[
\|C\| = \|m\|_{L^{\infty}(\mathbb{R})}.
\]
In particular, if $m(\omega)=a_{0}+\sum_{p}a_{p}\cos(\omega\log p)$
with $\{a_{p}\}\in\ell^{1}$, then
\[
\|C\| \leq |a_{0}| + \sum_{p}|a_{p}|.
\]
\end{proposition}

\begin{proof}
The Fourier multiplier $C=\mathcal{F}^{-1}M_{m}\mathcal{F}$ is
unitarily equivalent to multiplication by $m$ on
$L^{2}(\mathbb{R})$, so $\|C\|=\|M_{m}\|=\|m\|_{L^{\infty}}$.
The bound $\|m\|_{\infty}\leq|a_{0}|+\sum_{p}|a_{p}|$ is
immediate from the triangle inequality applied to the
Fourier series.
\end{proof}

\begin{theorem}[Global operator norm bound]
\label{thm:normU}
Under hypotheses \textup{(H1)--(H3)},
\[
\|U\| \;\leq\; \|A\| + \|C\| + \|E\|
\;\leq\; 2^{-\sigma} + \|K\|_{\mathrm{HS}}
       + \|m\|_{L^{\infty}} + \|E\|.
\]
\end{theorem}

\begin{proof}
For operators of the form
$T_{1}\otimes I\otimes I + I\otimes T_{2}\otimes I + I\otimes I\otimes T_{3}$
on a tensor product of Hilbert spaces,
$\|T_{1}\otimes I\otimes I\|=\|T_{1}\|$ (and analogously for the other terms),
so the triangle inequality gives
$\|U\|\leq\|A\|+\|C\|+\|E\|$.
The individual bounds come from
Propositions~\ref{prop:normA}--\ref{prop:normC}
and self-adjointness of $E$.
\end{proof}

% ------------------------------------------------------------
\subsection{Convergence and Boundedness of \texorpdfstring{$\Xi(t)$}{Xi(t)}}
\label{sec:Xi}
% ------------------------------------------------------------

\begin{theorem}[Convergence of recursive operator]
\label{thm:Xiconv}
Under hypothesis \textup{(H4)},
$\Xi(t)$ converges in the strong operator topology
for every $t\geq 0$, and
\[
\|\Xi(t)\| \;\leq\; \sum_{k=0}^{\infty}|\alpha_{k}(t)| \;\leq\; M.
\]
The family $\{\Xi(t)\}_{t\geq 0}$ is uniformly bounded by $M$.
\end{theorem}

\begin{proof}
Since $T$ is a bounded shift, $\|T^{k}\|\leq 1$ for all $k$.
By (H4),
\[
\left\|\sum_{k=0}^{N}\alpha_{k}(t)T^{k}\right\|
\leq \sum_{k=0}^{N}|\alpha_{k}(t)|\cdot\|T^{k}\|
\leq \sum_{k=0}^{N}|\alpha_{k}(t)|
\leq M.
\]
The partial sums form a Cauchy sequence in operator norm
by absolute convergence, so the series converges in norm.
Uniformity in $t$ follows from the uniform bound $M$.
\end{proof}

\begin{corollary}[Contraction semigroup criterion]
\label{cor:contract}
If additionally $M < 1$, then $\Xi(t)$ is a strict contraction
and $A:= -U$ generates a uniformly continuous contraction semigroup
$e^{tA}$ on $H$ provided $E\geq 0$ and $m(\omega)\geq 0$ a.e.
\end{corollary}

\begin{proof}
Under $E\geq 0$ (i.e.\ $\Xi\geq 0$ on $\mathbb{C}^{d}$) and
$m(\omega)\geq 0$, all three blocks of $U$ are positive
semi-definite, so $U\geq 0$.
Hence $A=-U\leq 0$ is a bounded dissipative operator.
By the Lumer--Phillips theorem (bounded version),
$A$ generates a uniformly continuous contraction semigroup,
and $\|e^{tA}\|\leq e^{-\lambda_{\min}(U)t}\leq 1$.
\end{proof}

% ------------------------------------------------------------
\subsection{Contractivity Index and the Correctness Inversion Lemma}
\label{sec:qt}
% ------------------------------------------------------------

\begin{definition}[Certification predicate]
\label{def:cert}
A state $\psi\in H$ \emph{passes} the contractivity certificate
with margin $\varepsilon>0$ if
\[
q_{t} \;=\; \|\Xi_{t}\| + \|\Lambda_{t}\|\,L_{T} \;<\; 1-\varepsilon.
\]
\end{definition}

\begin{lemma}[Correctness inversion under proxy fidelity]
\label{lem:inversion}
Let $\hat{q}$ be a \emph{proxy} fidelity score satisfying
$\hat{q}\in[0,1]$ but computed \emph{without} evaluating
$\|\Xi_{t}\|$ or $\|\Lambda_{t}\|$.
Then there exist operators $(\Xi_{t},\Lambda_{t})$ such that:
\begin{enumerate}[label=(\roman*)]
\item $\hat{q}\geq 0.95$ (proxy passes), yet $q_{t}\geq 1$ (system is unstable); and
\item $\hat{q}<0.95$ (proxy fails), yet $q_{t}<1-\varepsilon$ (system is stable).
\end{enumerate}
In particular, no threshold on $\hat{q}$ alone can certify
or refute contractivity.
\end{lemma}

\begin{proof}
For (i): Take $\Xi_{t}=c I$ with $c>1$ and $\Lambda_{t}=0$.
Then $q_{t}=c>1$.
Choose $\hat{q}=0.98$ as a hardcoded constant
independent of $(\Xi_{t},\Lambda_{t})$.
Both conditions hold simultaneously.

For (ii): Take $\Xi_{t}=\varepsilon' I$ with $0<\varepsilon'<1-\varepsilon$
and $\Lambda_{t}=0$, so $q_{t}=\varepsilon'<1-\varepsilon$.
Again set $\hat{q}=0.94<0.95$ independently.
Both conditions hold simultaneously.

Since in the current implementation
$\hat{q}$ is a static constant per marshal path
rather than a function of the actual operators,
both failure modes are realizable in production.
\end{proof}

\begin{remark}
Lemma~\ref{lem:inversion} formalises the ``silent correctness inversion''
identified in the body of the report.
It implies that repairing the proxy requires making
$\hat{q}$ a function of the actual operator inputs,
e.g.\ $\hat{q}\leftarrow f(\|\Xi_{t}\|,\|\Lambda_{t}\|,L_{T})$.
\end{remark}

% ------------------------------------------------------------
\subsection{Corrected Certification Protocol}
\label{sec:certprotocol}
% ------------------------------------------------------------

\begin{definition}[Corrected fidelity report]
\label{def:correctedfidelity}
Given inputs $(\Xi_{t},\Lambda_{t},L_{T},\varepsilon_{\min})$,
define the \emph{corrected fidelity report}
\[
\mathrm{FidelityReport}_{\mathrm{corr}}
= \left(\,\hat{q}_{\mathrm{corr}},\;
         q_{t},\;
         \mathrm{GapLB},\;
         \mathrm{SlopeUB}\,\right),
\]
where
\begin{align}
q_{t} &= \|\Xi_{t}\| + \|\Lambda_{t}\|\,L_{T}, \\
\hat{q}_{\mathrm{corr}} &= \max(0,\;1 - q_{t}), \\
\mathrm{GapLB}(w) &= \inf_{\theta}\delta_{S}(\theta) - 2\sum_{p}|w_{p}|b_{p}, \\
\mathrm{SlopeUB}(w) &= \sum_{p}|w_{p}|L_{p}.
\end{align}
The report is \emph{acceptable} if and only if
$q_{t} < 1-\varepsilon_{\min}$ (equivalently, $\hat{q}_{\mathrm{corr}}>\varepsilon_{\min}$).
\end{definition}

\begin{proposition}[Corrected acceptability is sufficient]
\label{prop:sufficiency}
If $\mathrm{FidelityReport}_{\mathrm{corr}}$ is acceptable,
then the sequence of states $\{\psi_{t}\}$ produced by
$\psi_{t+1}=\Xi_{t}\psi_{t}+\Lambda_{t}g_{t}$ satisfies
\[
\limsup_{t\to\infty}\|\psi_{t}\| \;\leq\;
\frac{\sup_{t}\|\Lambda_{t}\|\,\sup_{t}\|g_{t}\|}{1-(1-\varepsilon_{\min})}
= \frac{\sup_{t}\|\Lambda_{t}\|\,\sup_{t}\|g_{t}\|}{\varepsilon_{\min}}.
\]
In particular the sequence is bounded.
\end{proposition}

\begin{proof}
By induction,
\[
\|\psi_{t}\| \leq q_{t-1}\|\psi_{t-1}\| + \|\Lambda_{t-1}\|\,\|g_{t-1}\|
\leq (1-\varepsilon)\|\psi_{t-1}\| + B,
\]
where $B=\sup_{t}\|\Lambda_{t}\|\sup_{t}\|g_{t}\|$
and $1-\varepsilon = q_{t} < 1-\varepsilon_{\min}$.
Summing the geometric series gives the stated bound.
\end{proof}

% ------------------------------------------------------------
\subsection{Policy-to-Recurrence Interface Analysis}
\label{sec:policy}
% ------------------------------------------------------------

\begin{definition}[PIRTM policy protocol]
\label{def:protocol}
A \emph{PIRTM policy} is a triple of operator-valued maps
$(t\mapsto\Xi_{t},\;t\mapsto\Lambda_{t},\;t\mapsto G_{t})$
together with a scalar $\delta_{\mathrm{goal}}>0$ such that:
\begin{align}
\Xi_{t}&\in\mathcal{B}(V), \quad\text{for a Banach space }V, \\
\Lambda_{t}&\in\mathcal{B}(V,V), \\
G_{t}&\in V.
\end{align}
\end{definition}

\begin{lemma}[Carry-forward is not a PIRTM policy]
\label{lem:disconnect}
The concrete carry-forward hierarchy
$(\beta,\lambda)\mapsto\mathrm{CarryForwardPolicy}$
does not constitute a PIRTM policy in the sense of
Definition~\ref{def:protocol}.
\end{lemma}

\begin{proof}
By inspection of the implementation,
\texttt{CarryForwardPolicy} exposes:
$\texttt{beta}\in(0,10]$,
$\texttt{lambda\_decay}\in(0,1]$,
$\texttt{get\_next\_surplus}(S,\delta)\to\mathbb{R}$,
$\texttt{validate\_parameters}()\to\texttt{bool}$.
None of these is a map $t\mapsto\Xi_{t}$
from a time index to a bounded operator on a Banach space.
The \texttt{getattr} fallback in \texttt{iterate()} therefore
silently substitutes the default identity operator for $\Xi_{t}$
whenever a carry-forward policy is passed,
rendering the policy parameters inert.
\end{proof}

\begin{proposition}[Canonical adapter embedding]
\label{prop:adapter}
Define the adapter
\[
\Xi_{t}^{\mathrm{cfp}} = \lambda\,I_{V},
\qquad
\Lambda_{t}^{\mathrm{cfp}} = \frac{\beta}{10}\,I_{V},
\qquad
G_{t}^{\mathrm{cfp}} = 0,
\qquad
\delta_{\mathrm{goal}} = \varepsilon_{\mathrm{policy}},
\]
where $V=\mathbb{R}^{d}$.
Then:
\begin{enumerate}[label=(\roman*)]
\item $(\Xi_{t}^{\mathrm{cfp}},\Lambda_{t}^{\mathrm{cfp}},G_{t}^{\mathrm{cfp}})$
      is a valid PIRTM policy.
\item The contractivity index satisfies
      $q_{t}^{\mathrm{cfp}} = \lambda + \tfrac{\beta}{10}L_{T}$.
\item The carry-forward configuration satisfies
      $q_{t}^{\mathrm{cfp}}<1-\varepsilon$ if and only if
      \[
      \lambda + \tfrac{\beta}{10}L_{T} < 1-\varepsilon.
      \]
\end{enumerate}
\end{proposition}

\begin{proof}
(i): The maps $t\mapsto\lambda I$ and $t\mapsto\tfrac{\beta}{10}I$
are trivially bounded and time-independent, hence valid.

(ii): $\|\Xi_{t}^{\mathrm{cfp}}\|=\lambda$ and
$\|\Lambda_{t}^{\mathrm{cfp}}\|=\tfrac{\beta}{10}$,
so $q_{t}=\lambda+\tfrac{\beta}{10}L_{T}$ by Definition~\ref{def:qt}.

(iii): Immediate from (ii) and Definition~\ref{def:cert}.
\end{proof}

\begin{corollary}[Carry-forward contractivity criterion]
\label{cor:cfpcrit}
A carry-forward configuration $(\beta,\lambda)$ produces
a contractive dynamics under the canonical adapter embedding if and only if
\[
\lambda \;<\; 1 - \varepsilon - \frac{\beta\,L_{T}}{10}.
\]
For the default configuration invariants
$\lambda\leq 1$, $\beta\leq 10$, this constrains
$L_{T} < \tfrac{10(1-\varepsilon-\lambda)}{\beta}$.
\end{corollary}

% ------------------------------------------------------------
\subsection{Cross-Module Spectral Stability}
\label{sec:spectral}
% ------------------------------------------------------------

\begin{definition}[Session coupling system]
\label{def:coupling}
Let $n$ modules $M_{1},\dots,M_{n}$ have local
contraction factors $\gamma_{i}\in(0,1)$ and
let $C\in\mathbb{R}^{n\times n}_{\geq 0}$
be the coupling matrix where $C_{ij}$ represents
the strength of influence of module $j$ on module $i$.
\end{definition}

\begin{theorem}[Cross-module stability criterion]
\label{thm:crossmod}
Under Definition~\ref{def:coupling},
the composed system is globally stable
(all trajectories bounded) if and only if
\[
\rho\!\left(C\,\mathrm{diag}(\gamma_{1},\dots,\gamma_{n})\right) < 1,
\]
where $\rho(\cdot)$ denotes the spectral radius.
\end{theorem}

\begin{proof}
Let $\mathbf{x}_{t}=(x_{1,t},\dots,x_{n,t})^{\top}$
denote the stacked state vector.
The composed recurrence is
\[
\mathbf{x}_{t+1} = C\,\mathrm{diag}(\gamma_1,\dots,\gamma_n)\,\mathbf{x}_{t}
                 + \mathbf{g}_{t},
\]
for bounded forcing $\mathbf{g}_{t}$.
By the standard linear stability theorem,
the homogeneous part converges to zero for all initial conditions
if and only if all eigenvalues of
$\Gamma := C\,\mathrm{diag}(\gamma_{1},\dots,\gamma_{n})$
lie strictly inside the unit disc, i.e.\ $\rho(\Gamma)<1$.
\end{proof}

\begin{remark}
Note that individual module certification ($\gamma_{i}<1$ for all $i$)
is necessary but not sufficient for
$\rho(C\,\mathrm{diag}(\gamma_{1},\dots,\gamma_{n}))<1$.
For example, with $n=2$, $\gamma_{1}=\gamma_{2}=0.9$, and
$C=\begin{pmatrix}0&2\\2&0\end{pmatrix}$,
we get $\rho(\Gamma)=1.8>1$, so the composed system is unstable
despite each module being individually contractive.
This justifies the need for a session-level spectral check that is
not reducible to per-transition fidelity reports.
\end{remark}

\begin{proposition}[Sufficient condition via row sums]
\label{prop:rowsum}
If $C\geq 0$ and
$\sum_{j=1}^{n}C_{ij}\gamma_{j}<1$ for all $i$,
then $\rho(C\,\mathrm{diag}(\gamma_{1},\dots,\gamma_{n}))<1$.
\end{proposition}

\begin{proof}
Let $r_{i}=\sum_{j}C_{ij}\gamma_{j}<1$.
Then the matrix $\Gamma=C\,\mathrm{diag}(\gamma_{j})$
satisfies $\|\Gamma\|_{\infty}=\max_{i}r_{i}<1$.
Since $\rho(\Gamma)\leq\|\Gamma\|_{\infty}$, the result follows.
\end{proof}

\begin{corollary}[Runtime-checkable inequality]
\label{cor:runtimecheck}
The row-sum condition in Proposition~\ref{prop:rowsum}
is computable at session registration time given $C$ and $\{\gamma_{i}\}$,
providing an $O(n^{2})$ necessary check that can be enforced
without computing eigenvalues.
If the row-sum test passes, the session is pre-certified.
If it fails, the eigenvalue computation must proceed.
\end{corollary}

% ------------------------------------------------------------
\subsection{Ledger Conservation and Dual-Ledger Consistency}
\label{sec:ledger}
% ------------------------------------------------------------

\begin{definition}[Ledger surplus sequence]
\label{def:ledger}
Given a session $\mathcal{S}$ with $N$ phase transitions,
let $s_{k}\in\mathbb{R}$ be the surplus recorded at step $k$.
Define the \emph{total ledger balance}
$\mathcal{L}(\mathcal{S}) = \sum_{k=1}^{N}s_{k}$.
\end{definition}

\begin{definition}[Fidelity-induced surplus]
\label{def:surplus}
For a transition with fidelity score $f_{k}\in[0,1]$
and phase direction $\phi_{k}\in\{+1,-1\}$
(positive for tensor expansion, negative for atomic collapse),
define
\[
s_{k} = (f_{k}-1) + 0.01\,\phi_{k}.
\]
\end{definition}

\begin{proposition}[Ledger conservation bound]
\label{prop:ledgercons}
Under Definition~\ref{def:surplus},
\[
|\mathcal{L}(\mathcal{S})| \;\leq\; N(1-f_{\min}) + 0.01N,
\]
where $f_{\min}=\min_{k}f_{k}$.
The conservation criterion $|\mathcal{L}(\mathcal{S})|\leq\delta$
holds for any $\delta\geq N(1-f_{\min}+0.01)$.
\end{proposition}

\begin{proof}
For each step, $|s_{k}|\leq|f_{k}-1|+0.01\leq(1-f_{\min})+0.01$.
Summing over $k=1,\dots,N$ gives the stated bound.
\end{proof}

\begin{remark}[Ledger conservation vs.\ contractivity]
\label{rem:ledger}
The ledger conservation criterion
$|\mathcal{L}(\mathcal{S})|\leq 0.1$
(as implemented in \texttt{\_verify\_ledger\_conservation()})
is a bookkeeping check on surplus accumulation.
It is semantically orthogonal to the contractivity criterion
$q_{t}<1-\varepsilon$.
A session can satisfy ledger conservation while violating
contractivity, and vice versa.
The two conditions must therefore be enforced independently.
\end{remark}

% ------------------------------------------------------------
\subsection{ZM Certification Bounds}
\label{sec:ZM}
% ------------------------------------------------------------

\begin{definition}[Budgeted perturbation family]
\label{def:budget}
Let $U_{0}$ be a bounded self-adjoint operator with
spectral gap $\delta_{S}>0$ around a target band $S$.
Consider perturbations
\[
U(\theta) = U_{0} + \sum_{p\in\mathbb{P}} w_{p}\,B_{p}(\theta),
\]
with $\|B_{p}(\theta)\|\leq b_{p}$ and
$\|\partial_{\theta}B_{p}(\theta)\|\leq L_{p}$
for all $\theta$, and with budget $\sum_{p}|w_{p}|\leq \tau$.
\end{definition}

\begin{theorem}[ZM gap and slope certification]
\label{thm:ZMcert}
Under Definition~\ref{def:budget}, for all admissible $w$:
\begin{align}
\mathrm{GapLB}(w) &\geq \inf_{\theta}\,\delta_{S}(\theta)
               - 2\sum_{p}|w_{p}|\,b_{p}, \label{eq:GapLB}\\
\mathrm{SlopeUB}(w) &\leq \sum_{p}|w_{p}|\,L_{p}. \label{eq:SlopeUB}
\end{align}
\end{theorem}

\begin{proof}
By Weyl's inequality, adding a self-adjoint perturbation $P$
shifts eigenvalues by at most $\|P\|$:
\[
\left|\lambda_{k}(U(\theta))-\lambda_{k}(U_{0})\right|
\leq \left\|\sum_{p}w_{p}B_{p}(\theta)\right\|
\leq \sum_{p}|w_{p}|\|B_{p}(\theta)\|
\leq \sum_{p}|w_{p}|b_{p}.
\]
Since both band edges may move inward, the gap shrinks by at most
$2\sum_{p}|w_{p}|b_{p}$, giving \eqref{eq:GapLB}.

For \eqref{eq:SlopeUB}, differentiate and apply the triangle inequality:
\[
\|\partial_{\theta}U(\theta)\|
\leq \sum_{p}|w_{p}|\|\partial_{\theta}B_{p}(\theta)\|
\leq \sum_{p}|w_{p}|L_{p},
\]
and use the standard eigenband Lipschitz bound
from the spectral theorem.
\end{proof}

\begin{corollary}[Certifiability condition]
\label{cor:certifiable}
The gap target $\gamma_{\min}$ is achievable if and only if
\[
\tau \;<\; \frac{\inf_{\theta}\delta_{S}(\theta) - \gamma_{\min}}{2\,\sup_{p}b_{p}}.
\]
\end{corollary}

\begin{proof}
Setting $\sum_{p}|w_{p}|b_{p}\leq\tau\sup_{p}b_{p}$,
the condition $\mathrm{GapLB}(w)\geq\gamma_{\min}$
becomes $\inf_{\theta}\delta_{S}(\theta)-2\tau\sup_{p}b_{p}\geq\gamma_{\min}$.
Rearranging gives the stated inequality.
\end{proof}

\begin{proposition}[Hilbert--Schmidt threshold for ZM kernel]
\label{prop:HS}
The kernel operator $K$ from Definition~\ref{def:U}
is Hilbert--Schmidt if and only if $\alpha>\tfrac{1}{2}$.
\end{proposition}

\begin{proof}
\[
\|K\|_{\mathrm{HS}}^{2}
= \sum_{p,q\in\mathbb{P}} p^{-2\alpha}q^{-2\alpha}h(\log p-\log q)^{2}
\leq \|h\|_{\infty}^{2}\!\left(\sum_{p}p^{-2\alpha}\right)^{\!2}.
\]
By comparison with $\sum_{p}p^{-s}$,
the series converges if and only if $2\alpha>1$, i.e.\ $\alpha>\tfrac{1}{2}$.
\end{proof}

% ------------------------------------------------------------
\subsection{Finite-Prime Truncation Error}
\label{sec:truncation}
% ------------------------------------------------------------

\begin{theorem}[Truncation error bound]
\label{thm:truncation}
Let $A^{(N)}$ be the truncation of $A=D_{\sigma}+K$
to the first $N$ primes.
Let $p_{N+1}$ denote the $(N{+}1)$-th prime.
Then
\[
\|A - A^{(N)}\| \;\leq\; p_{N+1}^{-\sigma}
+ \|h\|_{\infty}\!\left[\sum_{p>p_{N}}p^{-2\alpha}
\right]^{1/2}
\!\!\left[\sum_{q}q^{-2\alpha}\right]^{1/2}.
\]
In particular, $\|A-A^{(N)}\|\to 0$ as $N\to\infty$
whenever $\alpha>\tfrac{1}{2}$ and $\sigma>0$.
\end{theorem}

\begin{proof}
The truncation error decomposes as
$\|A-A^{(N)}\|\leq\|D_{\sigma}-D_{\sigma}^{(N)}\|+\|K-K^{(N)}\|$.
The diagonal term satisfies
$\|D_{\sigma}-D_{\sigma}^{(N)}\|=p_{N+1}^{-\sigma}\to 0$.
For the off-diagonal term, by the Hilbert--Schmidt bound,
\[
\|K-K^{(N)}\|
\leq\|K-K^{(N)}\|_{\mathrm{HS}}
=\|h\|_{\infty}\!\left(\sum_{p>p_{N}}p^{-2\alpha}\right)^{1/2}
\!\left(\sum_{q}q^{-2\alpha}\right)^{1/2},
\]
which tends to zero as $N\to\infty$ by absolute convergence
of $\sum_{p}p^{-2\alpha}$ for $\alpha>\frac{1}{2}$.
\end{proof}

% ------------------------------------------------------------
\subsection{Summary of Operator Bounds}
\label{sec:summary}
% ------------------------------------------------------------

The following table consolidates the operator norms derived in this appendix.

\begin{table}[h!]
\centering
\caption{Summary of derived operator norm bounds}
\begin{tabular}{p{0.22\linewidth} p{0.36\linewidth} p{0.32\linewidth}}
\toprule
\textbf{Operator} & \textbf{Bound} & \textbf{Condition} \\
\midrule
$D_{\sigma}$ on $\ell^{2}(\mathbb{P})$
& $\|D_{\sigma}\|=2^{-\sigma}$
& $\sigma>0$ \\[4pt]
$K$ on $\ell^{2}(\mathbb{P})$
& $\|K\|\leq\|K\|_{\mathrm{HS}}=\|h\|_{\infty}\!\left(\sum_{p}p^{-2\alpha}\right)$
& $\alpha>\tfrac{1}{2}$ \\[4pt]
$A=D_{\sigma}+K$
& $\|A\|\leq 2^{-\sigma}+\|K\|_{\mathrm{HS}}$
& $\sigma>0,\,\alpha>\tfrac{1}{2}$ \\[4pt]
$C=\mathcal{F}^{-1}M_{m}\mathcal{F}$
& $\|C\|=\|m\|_{L^{\infty}}$
& $m\in L^{\infty}$ \\[4pt]
$U=A+C+E$ (tensor lifts)
& $\|U\|\leq\|A\|+\|m\|_{L^{\infty}}+\|E\|$
& (H1)--(H3) \\[4pt]
$\Xi(t)$
& $\|\Xi(t)\|\leq M$
& (H4) \\[4pt]
$q_{t}$ (contractivity index)
& $q_{t}=\|\Xi_{t}\|+\|\Lambda_{t}\|L_{T}$
& Definition~\ref{def:qt} \\[4pt]
$A^{(N)}$ truncation error
& $\|A-A^{(N)}\|\to 0$
& $\alpha>\tfrac{1}{2},\,\sigma>0$ \\[4pt]
\bottomrule
\end{tabular}
\end{table}

\begin{remark}[Roadmap to implementation]
Every bound in this table is computable given the explicit parameters
$(\sigma,\alpha,h,m,d,N)$ and the operator norms $\|\Lambda_{t}\|,L_{T}$.
These are precisely the arguments that should be passed to a corrected
\texttt{certify\_state()} implementation in order to eliminate the
proxy-fidelity inversion described in Lemma~\ref{lem:inversion}.
\end{remark}

% ============================================================
% END MATHEMATICAL APPENDIX
% ============================================================
\nocite{*}
\bibliographystyle{unsrt}  
\bibliography{references}
\end{document}